%% kapa: exam=1 round=34 year=2023 period=2 form=A subject=accounting subject_ko=회계학 seq=1 number=41
%% source: Q-Net 감정평가사 기출문제 (cst00309) artlSeq=5211789 file=raw/1008/5211789/2023년도 제34회 감정평가사 1차시험 2교시.pdf
%% conversion: kordoc 4.24.51 -> tools/gen1.py (auto); review=proofread-sonnet; flags=none
\begin{question}{41}
재무제표 표시에 관한 일반사항으로 옳지 않은 것은?
\begin{choices}
  \item 서술형 정보는 당기 재무제표를 이해하는 데 목적적합하더라도 비교정보를 표시하지 아니한다.
  \item 재무제표가 계속기업 기준으로 작성되지 않을 경우, 그 사실과 함께 재무제표 작성기준과 계속기업으로 보지 않는 이유를 공시하여야 한다.
  \item 기업은 현금흐름 정보를 제외하고는 발생기준 회계를 사용하여 재무제표를 작성한다.
  \item 중요하지 않은 항목은 성격이나 기능이 유사한 항목과 통합하여 표시할 수 있다.
  \item 한국채택국제회계기준을 준수하여 작성된 재무제표는 공정하게 표시된 재무제표로 본다.
\end{choices}
\end{question}
\answer{1}
